% New author continuation. No collision is inserted at a marked foot.
% Dependencies: 277, 279--283, 316--320, 336, and the angular lemma of 341.
\section{Exponential payment of the regular unvisited-boundary atoms}
\label{sec:umi-unvisited}

The previous section treated the original section sides. We now
pay the regular constant-roof pieces of every genuine unvisited
boundary. The source is exactly the marked-flight angular source
of Sections~\ref{sec:nm-kernel} and \ref{sec:mb-capacity}, including
its already proved positive-side zero-clearance trace. No limiting
weak kernel is used to infer a sum of branch suprema.

Fix $R\in[9/20,47/100]$, the true collision count $m\ge1$, an
actual flight index $0\le j<m$, and a prescribed clearance level
$u_0\in[0,3/200]$. At a positive level specify one incident side;
if both are counted their charges are added. At $u_0=0$ retain
precisely the physical $u>0$ unvisited trace. Let $\varsigma=\pm1$
be its impact sign and $\vartheta$ its line direction. The marked
disk is translated to the origin only to name the same source.
The foot is
\[
 q_f=-\varsigma(R+u_0)v_\vartheta^\perp.
\]
The original record is recovered by $j+1$ actual hits backwards
and $m-j$ actual hits forwards from this point. These are the
existing inclusive endpoint counts of
Section~\ref{sec:bf-foot-split}; the original roof still consists
of exactly $m$ physical flights.

Use the constants
\begin{equation}\label{eq:umi-constants}
 \begin{gathered}
 E_* =\frac{94}{\sqrt{1827}},\qquad
 C_{\rm foot}=\frac{607}{90},\qquad H=H_{\rm nm},\\
 Q=\frac{U}{c_\circ\kappa_-},\qquad
 C_{\rm umi}=F_{\rm cs}C_{\rm foot}HE_*(\pi+2Q).
 \end{gathered}
\end{equation}
Here $F_{\rm cs}$ and $Q$ are exactly
\eqref{eq:iso-constants}. In particular $H$ includes both impact
signs. It is the constant for the marked-foot curve class of
Theorem~\ref{thm:mb-branch-sup}; no constant belonging only to a
different source class is substituted.

\subsection{The marked source has the same true endpoint balance}

\begin{lemma}[Stationary roof on an actual marked angular source]
\label{lem:umi-balance}
On each regular marked overlap, including the designated
positive-visited zero-clearance trace, the actual endpoint
momenta obey
\begin{equation}\label{eq:umi-balance}
 \partial_\vartheta F
   =\frac{p_0\partial_\vartheta p_0}{c_0V_-}
       +\frac{p_m\partial_\vartheta p_m}{c_mV_+},
 \qquad \kappa_-\le V_-,V_+\le U.
\end{equation}
On the original endpoint cap, every regular stationary-roof
subsource therefore has the single-opposition-quadrant and
thin-angular-wedge conclusions of Lemma~\ref{lem:iso-balance},
with parameter $\vartheta$.
\end{lemma}
\begin{proof}
At fixed $u_0$ the derivative of $q_f$ is parallel to
$v_\vartheta$. Its transverse field is $(0,1)$, in both time
directions with the inherited reversal convention. By
Theorem~\ref{thm:ag-density}, each first true image is an unstable
graph with slope between $\kappa_-$ and
$\kappa_++a_\circ^{-1}<U$. The actual graph transform preserves
$[\kappa_-,U]$. Thus the two endpoint graph identities are
\[
 \partial_\vartheta p_0
      =-Rc_0V_-\partial_\vartheta\alpha_0,
 \qquad
 \partial_\vartheta p_m
      =Rc_mV_+\partial_\vartheta\alpha_m.
\]
Insert them into the true endpoint first variation
$dF=-Rp_0\,d\alpha_0+Rp_m\,d\alpha_m$.
If the physical path is split at $q_f$ to perform this calculation,
the two contributions $v_\vartheta\cdot dq_f$ and
$-v_\vartheta\cdot dq_f$ cancel exactly. The point is not a
reflection or an additional event. This proves
\eqref{eq:umi-balance} with the original physical clock.

Both endpoint momentum derivatives are nonzero with fixed signs
on each actual half-word overlap, by
\eqref{eq:bf-signed-derivatives}. At a stationary point away from
$(p_0,p_m)=(0,0)$, neither component can vanish. After orienting
both components increasingly, stationarity places them in one
opposition quadrant. The proof of Lemma~\ref{lem:iso-balance}
then applies verbatim: the Gauss angle is one-to-one on that
whole quadrant interval. This concerns the entire selected
subset, not only each of its connected pieces. The momentum
ratio bound is $|p_{\rm long}|/|p_{\rm short}|\le Qb/a$ for
the corresponding inverse endpoint derivatives $a,b$.

At $u_0=0$, Theorem~\ref{thm:nm-trace} supplies the same
positive-visited analytic words and first-image formulas directly
on the specified one-sided source. The marked disk can be hit
later as an ordinary scatterer; only its designated foot tangency
remains unvisited. The proof introduces no derivative of a
collision map at that tangency.
\end{proof}

Let $\mathfrak A^{\rm umi}_{j,u_0}$ denote the complete one-copy
regular charge $A|d\theta|$ on this marked source restricted to
$\partial_\vartheta F=0$ and $(p_0,p_m)\ne(0,0)$. Both impact
signs and all actual marked copies are included. Original exact
record, physical roof, section and positive guard restrictions
can be retained. The restriction is used only to estimate
already defined physical derivative terms; it does not define
the derivative of an arbitrary new Borel event.

\begin{theorem}[All-flight marked isoroof payment]
\label{thm:umi-charge}
For every permitted $R,m,j,u_0$,
\begin{equation}\label{eq:umi-charge}
 \mathfrak A^{\rm umi}_{j,u_0}
 \le C_{\rm umi}\,
           \sigma^{\max(j,m-1-j)/2}.
\end{equation}
Consequently, for any prescribed levels
$u_j\in[0,3/200]$ with one incident side specified at each
positive level,
\begin{equation}\label{eq:umi-all-flights}
 \sum_{j=0}^{m-1}\mathfrak A^{\rm umi}_{j,u_j}
 \le B_{\rm umi}(m):=C_{\rm umi}m\sigma^{(m-1)/4}.
\end{equation}
All regular constant-roof pieces of the physical unvisited
boundary are included by taking $u_j=0$.
\end{theorem}
\begin{proof}
The actual half-counts are
$l_-=j+1$ and $l_+=m-j$. Designate the larger half-count as
the long half and break equality by a fixed rule. Put
\[
 L=\max(j+1,m-j),\qquad s=\min(j+1,m-j),\qquad
 \epsilon=E_*\sigma^{L-1},\quad A_s=E_*\sigma^{s-1}.
\]
Theorem~\ref{thm:mb-branch-sup} gives actual disjoint angular
interval families, on the same two labelled circles, with
\[
 \sum_Ia_I\le H,\quad \sum_Jb_J\le H,\qquad
 a_I\le A_s\le E_*,\quad b_J\le\epsilon\le E_*.
\]
The weights are the full interval suprema of the inverse short
and long endpoint momentum derivatives. Endpoint cap restrictions
leave at most one interval in each half-source interval. These
same branch-supremum statements hold directly at $u_0=0$ by the
final paragraph of that theorem's proof; no $u\downarrow0$ limit
is interchanged with a sum of suprema.

The exact marked determinant and physical coarea give
\[
 u_x=\frac1{R A_-^p A_+^p(h_-+h_+)}
       \le C_{\rm foot}a(\vartheta)b(\vartheta),
 \qquad A\le F_{\rm cs}u_x.
\]
The first inequality already contains the radius factor from
the marked line coordinates. No additional $R^{-1}$ or
$c_*^{-1}$ is inserted here.

For a cutoff $\eta>0$, on one overlap $I\cap J$ the complete
stationary subset with $a\le\eta$ has charge at most
\[
 F_{\rm cs}C_{\rm foot}\min(a_I,\eta)b_J\,\pi/2.
\]
Its complementary subset has charge at most
\[
 F_{\rm cs}C_{\rm foot}a_Ib_J\,Q\epsilon/\eta.
\]
Indeed Lemma~\ref{lem:umi-balance} places the first subset in
one monotone quadrant and the second in one monotone wedge of
width at most $Q\epsilon/\eta$. Disconnected selected pieces
are not counted repeatedly.

Apply Lemma~\ref{lem:mb-overlay} to the original interval
families, first with weights $\min(a_I,\eta),b_J$, then with
$a_I,b_J$. The two sums are at most
$H(\epsilon+\eta)$ and $H(\epsilon+A_s)$, respectively.
It follows that
\[
 \mathfrak A^{\rm umi}_{j,u_0}
 \le F_{\rm cs}C_{\rm foot}H
 \left\{\frac\pi2(\epsilon+\eta)
           +\frac{Q\epsilon(\epsilon+A_s)}\eta\right\}.
\]
Take $\eta=\sqrt{E_*\epsilon}$. As in the explicit calculation
in Theorem~\ref{thm:iso-interior}, the braces are at most
$(\pi+2Q)\sqrt{E_*\epsilon}$.
Since $L-1=\max(j,m-1-j)$, this proves
\eqref{eq:umi-charge}, including $j=0$, $j=m-1$ and $m=1$.
There is no second count of the two impact signs: both are
already in $H$ and in the disjoint labelled-source overlay.

The physical uniqueness argument in
Theorem~\ref{thm:mb-two-sided-sum} identifies each angular point
and sign with its original marked copy at the prescribed flight.
Recentring the initial disk creates no lattice multiplicity.
If several marks describe one physical interface, assign that
interface once by a fixed ordering and restrict the positive
marked sum. No bound on the number of its descriptions is
multiplied into the estimate.

All regular pieces and their one-sided trace conventions are
exactly those of Theorem~\ref{thm:ab-all-offset}. Compact regular
subarcs and finite positive families exhaust this complete
source. Additional isolated singular angular parameters have
zero regular angular charge on positive-visited noncritical
charts, where $A$ and the Gauss map are locally smooth.
Visited-grazing limits retain the separately proved zero trace;
critical corners retain their incoming sectors and belong to
Theorem~\ref{pc:fan}. No angular-null statement is substituted
for that critical trace theorem.

Finally $\max(j,m-1-j)\ge(m-1)/2$. Summing the already complete
per-flight bound gives \eqref{eq:umi-all-flights}. This finite
uniform inequality precedes every radius, label and roof
supremum and every subsequent collision limit.
\end{proof}

\begin{corollary}[The actual signed unvisited regular atoms]
\label{cor:umi-signed}
Let $b^{\rm umi}_\ell(t)$ be the signed contributions in
\eqref{eq:sbv-atomic-coefficient} from regular constant-roof
unvisited interfaces, assigned once and combined at the same
exact label and physical roof. For the original admitted
observations and guard,
\begin{equation}\label{eq:umi-signed}
 \sum_{\ell,t}|b^{\rm umi}_\ell(t)|
       \le2a_0b_0 B_{\rm umi}(m),\qquad
 \sup_R m^q\sum_{\ell,t}|b^{\rm umi}_\ell(t)|\longrightarrow0
 \quad(q\ge0\text{ fixed}).
\end{equation}
\end{corollary}
\begin{proof}
Use the complete source at $u_j=0$ and the original one-sided
record selectors. The difference of the two incident exact-label
vectors has total absolute norm at most two. This is a
conservative bound also for a physical exterior side, where
one vector is zero. Multiply by $|V|\le a_0b_0$ and
$0\le w\le1$. Combine signed terms at the same roof and label
before the triangle inequality. Apply
\eqref{eq:umi-all-flights}. The convergence follows from
$m^{q+1}\sigma^{(m-1)/4}\to0$; no outer parameter depends on
$m$ in this passage.
\end{proof}

\paragraph{Remaining derivative and the original target.}
Together with the preceding original-section results, this pays
all regular constant-roof physical boundary terms in
\eqref{eq:sbv-atomic-coefficient}. Critical germs on an unvisited
boundary are paid by the retained zero-clearance theorem, and
critical germs on an original section boundary by
Theorem~\ref{thm:scf-section-payment}. Artificial seams cancel
and visited-grazing limits have their retained zero trace.
The auxiliary endpoint-cap sides in the geometric classification
have zero exact-event label traces: both original endpoint
section tests lie strictly inside $|p|<5/6$, hence away from
the auxiliary cap $|p|=5/6$ and from the larger estimating
cap $|p|=9/10$.
Thus the only unpaid atomic source in that classification is
the positive-visited critical source with positive unvisited
clearance strictly inside its original section chamber.
The absolutely continuous derivative and positive-width source
remainders still require their original normalized estimates.
The compensated primitive of Proposition~\ref{prop:pa-primitive}
applies to the enlarged paid atomic collection with the sum of
these explicit bounds. This changes neither the physical density,
the exact microscopic event, nor the order of limits in the
original raw local-limit and inverse problem.
